\section{Controller Scope and Architecture}
\label{sec:architecture}

\subsection{Boundary of the Model}

UCIe separates protocol, die-to-die adapter, and physical-layer concerns, with
the sideband supporting training, register access, and diagnostics
\cite{sharma2022ucie,nature2024ucie3d}. This work addresses a digital-control
slice within that broader stack. As Fig.~\ref{fig:scope} shows, the environment
provides transaction-level sideband responses, training samples, and abstract
phase completions. The controller returns state/status, sideband requests,
training patterns, and retained diagnostics. Analog signaling, the package
channel, mainband data transfer, complete adapter/protocol operation, and
compliance are deliberately outside the boundary.

\begin{figure*}[t]
\centering
\includegraphics[width=0.97\textwidth]{fig01_scope.pdf}
\caption{Implemented RTL and evidence boundary. Colored blocks identify the
synthesized modules and their test/observer interfaces; the pale dashed panel
lists excluded physical-layer and full-protocol functions.}
\label{fig:scope}
\end{figure*}

The modeled state and substate names follow UCIe Specification Revision 2.0,
Version 1.0 \cite{ucie20}. Most MBINIT, MBTRAIN, and PHYRETRAIN operations are
represented by an external \texttt{phase\_done\_i} condition. One training
operation, DATATRAINCENTER1, is implemented through the sideband and LFSR
engines. The distinction prevents a control-flow model from being mistaken for
a complete PHY.

\subsection{Module Decomposition}

Fig.~\ref{fig:rtl} shows the synthesizable hierarchy. The top-level
\texttt{ucie\_ltsm} owns state/substate registers, a residency timer, routing,
and engine coordination. Three engines isolate operations whose temporal
contracts differ from ordinary state progression.

\begin{figure*}[t]
\centering
\includegraphics[width=0.96\textwidth]{fig02_rtl_architecture.pdf}
\caption{Synthesizable hierarchy. The controller instantiates the sideband,
LFSR, and error engines. The separate FPGA top instantiates the same controller
and serializes its wide diagnostics through a compact CSR interface.}
\label{fig:rtl}
\end{figure*}

The sideband sequencer has IDLE, SEND, WAIT, and ERROR states. It holds the
request under ready/valid backpressure, accepts only the expected response,
retries a bounded number of response timeouts, and exposes protocol error and
abort status. The current message enumeration includes NOP, SBINIT-DONE
request/response, and DATACENTER1 START/END request/response. It is not a
complete sideband packet layer.

The LFSR engine maintains 16 independent 23-bit states. Eight nonzero seeds are
repeated over the lanes, and the feedback polynomial is
$x^{23}+x^{21}+x^{16}+x^8+x^5+x^2+1$. State and sample count advance only when
\texttt{train\_rx\_valid\_i} accepts a sample. Mismatches accumulate in a
16-bit saturating counter, and pass is asserted only when the final count is
strictly less than the programmed threshold; equality fails. Fig.~\ref{fig:lfsr}
summarizes the datapath and its independent reference boundary.

\begin{figure*}[t]
\centering
\includegraphics[width=0.86\textwidth]{fig05_lfsr_engine.pdf}
\caption{LFSR measurement engine. Receive gaps freeze both the sample count and
per-lane expected state.}
\label{fig:lfsr}
\end{figure*}

Fig.~\ref{fig:wave-lfsr} is a VCD-backed module-directed capture rather than a
hand-drawn timing sketch. The receive-valid gaps leave the expected lane state
and error count unchanged; the accepted TX/RX samples match, so the engine
asserts done and pass after the eighth accepted sample. Signal names in the
standalone SVG link to the repository's public
\href{https://github.com/Varadain/ucie-ltsm/blob/feature/v1.0-integrated-training/docs/02_ltssm/signals.md}{signal-definition guide}.

\begin{figure*}[t]
\centering
\includegraphics[width=0.97\textwidth]{fig09_wave_lfsr.pdf}
\caption{Color-coded Questa waveform generated from
\texttt{build/waves/datatrain\_lfsr.vcd}. The capture shows sample-gated LFSR
progression, an exact zero mismatch count, and the final done/pass decision.}
\label{fig:wave-lfsr}
\end{figure*}

The error manager accepts one eligible state-timeout, sideband-protocol, or
local-fatal event. Simultaneous events use the fixed priority timeout,
sideband, then fatal. A pending event holds the TRAINERROR request until an
acknowledgment or manager timeout permits entry. Cause and a saturating 16-bit
event count remain stable in TRAINERROR. A software clear is accepted only
when no event is pending and the controller is outside TRAINERROR.

\subsection{Hierarchical LTSM Control}

The package enumerates nine top states: RESET, SBINIT, MBINIT, MBTRAIN,
LINKINIT, ACTIVE, PHYRETRAIN, TRAINERROR, and L1/L2. MBINIT contains six
substates and MBTRAIN contains thirteen. The nominal route progresses through
initialization and training to ACTIVE. From ACTIVE, requests can select one of
three retraining targets or enter the modeled low-power state; exit conditions
return to ACTIVE. Eligible timeout, protocol, or fatal conditions route through
the retained error manager to TRAINERROR, as shown in
Fig.~\ref{fig:ltsmflow}.

\begin{figure*}[t]
\centering
\includegraphics[width=0.93\textwidth]{fig03_ltsm_flow.pdf}
\caption{Hierarchical controller flow. Only DATATRAINCENTER1 is expanded here;
other physical operations remain phase-completion abstractions.}
\label{fig:ltsmflow}
\end{figure*}

One timer measures state/substate residency. It restarts on a state or substate
change and can be stalled by the environment. Default parameters correspond to
an 80-MHz controller clock, a 4-ms reset interval, an 8-ms residency timeout,
a 256-cycle sideband response timeout, one retry, and 4096 training samples.
Testbenches scale selected values for runtime while retaining the ordering
properties under test.

\subsection{Integrated DATATRAINCENTER1 Protocol}

Production mode fixes the abstract CENTER1 bypass to zero. On entering
DATATRAINCENTER1, the controller orders the sideband engine to send START and
wait for its matching response. Only then may the LFSR engine begin. A failed
pattern re-enters the pattern phase without emitting END; a passing pattern
starts the END exchange. Advancement to DATATRAINVREF requires the matching END
response. A mismatch, exhausted sideband retry, fatal event, state timeout, or
reset aborts the sequence along its corresponding control path.

\begin{figure*}[t]
\centering
\includegraphics[width=0.93\textwidth]{fig04_integrated_sequence.pdf}
\caption{Integrated CENTER1 transaction order. The independent predictor checks
the same interface events without reading the DUT's private LFSR state.}
\label{fig:center1}
\end{figure*}

This ordering, illustrated in Fig.~\ref{fig:center1}, creates two explicit
commit barriers: the START response admits pattern activity, and the END
response admits state advancement. Separating the sideband and pattern engines
keeps retry time from advancing training state and keeps receive gaps from
consuming the sideband timeout counter.

\subsection{Retained Diagnostics and FPGA CSR}

The compact wrapper preserves the controller interface internally but replaces
wide state/error/training diagnostics with byte reads. Read addresses 0x00,
0x01, and 0x02 expose top, MBINIT, and MBTRAIN state; 0x03 exposes status bits;
0x04 exposes error cause; 0x05--0x06 and 0x07--0x08 expose retained-event and
training-error counts; 0x09 returns version 0x04; and 0x0A exposes the integrated
training phase. Writing bit 0 at 0x10 requests a protected log clear. Undefined
reads return zero.

Fig.~\ref{fig:recovery} connects the retained log, entry handshake, recovery
route, CSR view, and assertion boundary. This observability choice is not only
documentation convenience: Sec.~\ref{sec:evaluation} shows that exporting wide
diagnostics prevents the chosen FPGA from fitting, whereas the CSR wrapper
fits.

\begin{figure*}[t]
\centering
\includegraphics[width=0.86\textwidth]{fig06_error_recovery.pdf}
\caption{Retained error and recovery flow. Clear is ignored while an event is
pending or the state is TRAINERROR.}
\label{fig:recovery}
\end{figure*}

The VCD-backed capture in Fig.~\ref{fig:wave-error} makes the temporal contract
visible. A one-cycle fatal event is retained, the request remains asserted
while pending, and a clear attempt is ignored before the bounded TRAINERROR
entry. The retained cause and event count survive until recovery reaches an
eligible clear state.

\begin{figure*}[t]
\centering
\includegraphics[width=0.97\textwidth]{fig10_wave_trainerror.pdf}
\caption{Color-coded Questa waveform from the self-checking error-manager
testbench. Colored phase bands identify fault capture, protected pending time,
bounded TRAINERROR entry, and the allowed post-recovery clear.}
\label{fig:wave-error}
\end{figure*}

Fig.~\ref{fig:wave-csr} then exposes the same retained behavior through the
practical FPGA top. CSR reads observe state, cause, and count; a write to 0x10
cannot clear TRAINERROR, but the same control succeeds after reset. The final
invalid read returns zero as specified by the wrapper map.

\begin{figure*}[t]
\centering
\includegraphics[width=0.97\textwidth]{fig11_wave_csr.pdf}
\caption{Color-coded Questa waveform from the FPGA-wrapper directed test. The
CSR sequence reads retained diagnostics, demonstrates protected and allowed
clear operations, and checks an undefined address.}
\label{fig:wave-csr}
\end{figure*}
